% GENERATED FILE. Edit canonical lessons, then run npm run generate:books.
% canonical-source-hash: fnv1a64:b6cb0776d325ab0b
% canonical-lessons: TE-C412-parisodhana, TE-C412-parisodhincu, TE-C412-parisodhakudu, TE-C412-samacaram, TE-C412-anvesincu

\chapter{Research}
\label{ch:te-research}
\hlchaptermodality{412}

\begin{chapteropening}
\textbf{By the end of this chapter:} \emph{I can talk about research, information and exploring a question.}

Complete the last lesson of chapter 412: I can talk about research, information and exploring a question.
\end{chapteropening}

\section[\texorpdfstring{\te{పరిశోధన} (pariśōdhana)}{పరిశోధన (pariśōdhana)}]{\te{పరిశోధన} (pariśōdhana) — research}

\label{lesson:TE-C412-parisodhana}



\begin{quote}
Before the new one: say the Telugu for public opinion, then the Telugu for a percentage.
\end{quote}

\subsection*{You'll want to know: \te{పరిశోధన}}
\textbf{\te{పరిశోధన}} — \emph{pariśōdhana} — \textquotedblleft{}research\textquotedblright{}.

\begin{quote}
\te{ఆమె} \te{తెలుగు} \te{పాటల} \te{మీద} \te{పరిశోధన} \te{చేస్తోంది}.
\end{quote}

\begin{quote}
\emph{āme telugu pāṭala mīda pariśōdhana cēstōndi.}
\end{quote}

\begin{quote}
She is doing research on Telugu songs.
\end{quote}

\te{పరిశోధన} \te{చేయు}, to do research.

\begin{cousinweb}
\textbf{\te{పరిశోధన}} (\emph{pariśōdhana}) is borrowed from Sanskrit \emph{pariśodhana}, \textquotedblleft{}cleaning, purification\textquotedblright{} (Monier-Williams); in Telugu it is going through a subject thoroughly, research.
\end{cousinweb}

\begin{grammarlens}[title={What you've built}]
One more word for giving reasons and opinions.
\end{grammarlens}

\subsection*{Your turn}
\begin{itemize}
\raggedright
  \item \emph{Say it:} \emph{pariśōdhana}
  \item \emph{Say it:} \emph{pariśōdhana}, once more
  \item \emph{Say it:} \emph{āme telugu pāṭala mīda pariśōdhana cēstōndi.} — She is doing research on Telugu songs.
  \item \emph{Recall:} say the Telugu for public opinion, then the Telugu for a percentage, then say \emph{pariśōdhana} again
\end{itemize}

\begin{tcolorbox}[breakable,colback=teal!4,colframe=teal!35!black,title={Before you move on}]
What does \te{పరిశోధన} mean? (\textquotedblleft{}Research\textquotedblright{}.) Say it once more.
\end{tcolorbox}
\section[\texorpdfstring{\te{పరిశోధించు} (pariśōdhiñcu)}{పరిశోధించు (pariśōdhiñcu)}]{\te{పరిశోధించు} (pariśōdhiñcu) — to research; to investigate}

\label{lesson:TE-C412-parisodhincu}



\begin{quote}
Before the new one: say the Telugu for a percentage, then the Telugu for research.
\end{quote}

\subsection*{You'll want to know: \te{పరిశోధించు}}
\textbf{\te{పరిశోధించు}} — \emph{pariśōdhiñcu} — \textquotedblleft{}to research; to investigate\textquotedblright{}.

\begin{quote}
\te{ఆమె} \te{తన} \te{కుటుంబ} \te{చరిత్రను} \te{పరిశోధించింది}.
\end{quote}

\begin{quote}
\emph{āme tana kuṭumba caritranu pariśōdhiñcindi.}
\end{quote}

\begin{quote}
She researched her family history.
\end{quote}

Past (I …): \te{పరిశోధించాను}, \emph{pariśōdhiñcānu}.

Also for the police: a case is investigated.

\begin{cousinweb}
\textbf{\te{పరిశోధించు}} (\emph{pariśōdhiñcu}) is built on Sanskrit \emph{pariśodhana}, \textquotedblleft{}cleaning, purification\textquotedblright{} (Monier-Williams), with the Telugu verb ending -\te{ఇంచు}; the noun is \te{పరిశోధన}.
\end{cousinweb}

\begin{grammarlens}[title={What you've built}]
One more word for giving reasons and opinions.
\end{grammarlens}

\subsection*{Your turn}
\begin{itemize}
\raggedright
  \item \emph{Say it:} \emph{pariśōdhiñcu}
  \item \emph{Say it:} \emph{pariśōdhiñcu}, once more
  \item \emph{Say it:} \emph{āme tana kuṭumba caritranu pariśōdhiñcindi.} — She researched her family history.
  \item \emph{Recall:} say the Telugu for a percentage, then the Telugu for research, then say \emph{pariśōdhiñcu} again
\end{itemize}

\begin{tcolorbox}[breakable,colback=teal!4,colframe=teal!35!black,title={Before you move on}]
What does \te{పరిశోధించు} mean? (\textquotedblleft{}To research; to investigate\textquotedblright{}.) Say it once more.
\end{tcolorbox}
\section[\texorpdfstring{\te{పరిశోధకుడు} (pariśōdhakuḍu)}{పరిశోధకుడు (pariśōdhakuḍu)}]{\te{పరిశోధకుడు} (pariśōdhakuḍu) — a researcher}

\label{lesson:TE-C412-parisodhakudu}



\begin{quote}
Before the new one: say the Telugu for research, then the Telugu for to research.
\end{quote}

\subsection*{You'll want to know: \te{పరిశోధకుడు}}
\textbf{\te{పరిశోధకుడు}} — \emph{pariśōdhakuḍu} — \textquotedblleft{}a researcher\textquotedblright{}.

\begin{quote}
\te{ఆ} \te{పరిశోధకుడు} \te{పది} \te{సంవత్సరాలు} \te{అడవిలో} \te{ఉన్నాడు}.
\end{quote}

\begin{quote}
\emph{ā pariśōdhakuḍu padi saṁvatsarālu aḍavilō unnāḍu.}
\end{quote}

\begin{quote}
That researcher stayed in the forest for ten years.
\end{quote}

Plural: \te{పరిశోధకులు}, \emph{pariśōdhakulu}.

A woman researcher is \te{పరిశోధకురాలు}.

\begin{cousinweb}
\textbf{\te{పరిశోధకుడు}} (\emph{pariśōdhakuḍu}) is built on Sanskrit \emph{pariśodhana} (Monier-Williams), with -\te{అక}, \textquotedblleft{}one who\textquotedblright{}, and the Telugu masculine ending -\te{డు}.
\end{cousinweb}

\begin{grammarlens}[title={What you've built}]
One more word for giving reasons and opinions.
\end{grammarlens}

\subsection*{Your turn}
\begin{itemize}
\raggedright
  \item \emph{Say it:} \emph{pariśōdhakuḍu}
  \item \emph{Say it:} \emph{pariśōdhakuḍu}, once more
  \item \emph{Say it:} \emph{ā pariśōdhakuḍu padi saṁvatsarālu aḍavilō unnāḍu.} — That researcher stayed in the forest for ten years.
  \item \emph{Recall:} say the Telugu for research, then the Telugu for to research, then say \emph{pariśōdhakuḍu} again
\end{itemize}

\begin{tcolorbox}[breakable,colback=teal!4,colframe=teal!35!black,title={Before you move on}]
What does \te{పరిశోధకుడు} mean? (\textquotedblleft{}A researcher\textquotedblright{}.) Say it once more.
\end{tcolorbox}
\section[\texorpdfstring{\te{సమాచారం} (samācāraṁ)}{సమాచారం (samācāraṁ)}]{\te{సమాచారం} (samācāraṁ) — information}

\label{lesson:TE-C412-samacaram}



\begin{quote}
Before the new one: say the Telugu for to research, then the Telugu for a researcher.
\end{quote}

\subsection*{You'll want to know: \te{సమాచారం}}
\textbf{\te{సమాచారం}} — \emph{samācāraṁ} — \textquotedblleft{}information\textquotedblright{}.

\begin{quote}
\te{ఆ} \te{విషయం} \te{మీద} \te{నాకు} \te{ఎక్కువ} \te{సమాచారం} \te{కావాలి}.
\end{quote}

\begin{quote}
\emph{ā viṣayaṁ mīda nāku ekkuva samācāraṁ kāvāli.}
\end{quote}

\begin{quote}
I need more information on that subject.
\end{quote}

\te{సమాచారం} \te{ఇచ్చు}, to inform.

\begin{cousinweb}
\textbf{\te{సమాచారం}} (\emph{samācāraṁ}) is borrowed from Sanskrit \emph{samācāra}, \textquotedblleft{}conduct, practice, custom\textquotedblright{} (Monier-Williams); Telugu uses it for news and information.
\end{cousinweb}

\begin{grammarlens}[title={What you've built}]
One more word for giving reasons and opinions.
\end{grammarlens}

\subsection*{Your turn}
\begin{itemize}
\raggedright
  \item \emph{Say it:} \emph{samācāraṁ}
  \item \emph{Say it:} \emph{samācāraṁ}, once more
  \item \emph{Say it:} \emph{ā viṣayaṁ mīda nāku ekkuva samācāraṁ kāvāli.} — I need more information on that subject.
  \item \emph{Recall:} say the Telugu for to research, then the Telugu for a researcher, then say \emph{samācāraṁ} again
\end{itemize}

\begin{tcolorbox}[breakable,colback=teal!4,colframe=teal!35!black,title={Before you move on}]
What does \te{సమాచారం} mean? (\textquotedblleft{}Information\textquotedblright{}.) Say it once more.
\end{tcolorbox}
\section[\texorpdfstring{\te{అన్వేషించు} (anvēṣiñcu)}{అన్వేషించు (anvēṣiñcu)}]{\te{అన్వేషించు} (anvēṣiñcu) — to explore; to search for}

\label{lesson:TE-C412-anvesincu}



\begin{quote}
Before the new one: say the Telugu for a researcher, then the Telugu for information.
\end{quote}

\subsection*{You'll want to know: \te{అన్వేషించు}}
\textbf{\te{అన్వేషించు}} — \emph{anvēṣiñcu} — \textquotedblleft{}to explore; to search for\textquotedblright{}.

\begin{quote}
\te{వాళ్ళు} \te{నీళ్ళ} \te{కోసం} \te{కొత్త} \te{చోట్లు} \te{అన్వేషించారు}.
\end{quote}

\begin{quote}
\emph{vāḷḷu nīḷḷa kōsaṁ kotta cōṭlu anvēṣiñcāru.}
\end{quote}

\begin{quote}
They searched new places for water.
\end{quote}

Past (I …): \te{అన్వేషించాను}, \emph{anvēṣiñcānu}.

The noun is \te{అన్వేషణ}, a search.

\begin{cousinweb}
\textbf{\te{అన్వేషించు}} (\emph{anvēṣiñcu}) is built on Sanskrit \emph{anveṣaṇa}, \textquotedblleft{}seeking for, searching, investigating\textquotedblright{} (Monier-Williams), with the Telugu verb ending -\te{ఇంచు}.
\end{cousinweb}

\begin{grammarlens}[title={What you've built}]
One more word for giving reasons and opinions.
\end{grammarlens}

\subsection*{Your turn}
\begin{itemize}
\raggedright
  \item \emph{Say it:} \emph{anvēṣiñcu}
  \item \emph{Say it:} \emph{anvēṣiñcu}, once more
  \item \emph{Say it:} \emph{vāḷḷu nīḷḷa kōsaṁ kotta cōṭlu anvēṣiñcāru.} — They searched new places for water.
  \item \emph{Recall:} say the Telugu for a researcher, then the Telugu for information, then say \emph{anvēṣiñcu} again
\end{itemize}

\begin{tcolorbox}[breakable,colback=teal!4,colframe=teal!35!black,title={Before you move on}]
What does \te{అన్వేషించు} mean? (\textquotedblleft{}To explore; to search for\textquotedblright{}.) Say it once more.
\end{tcolorbox}
